\section{Methodology}\label{sec:method}

\subsection{Research questions and hypotheses}\label{sec:rqs}

The study asks four questions, obtained by splitting the two-part RQ2 of the proposal into
its threshold component and its comparison component. Each carries a directional hypothesis
taken from the approved proposal; both thresholds (80\,\% and 60\,\%) and the predicted sign
of the correlation were fixed there before any data were collected, so none of them is a
post-hoc expectation.

\textbf{RQ1 --- Coverage threshold.} When \texttt{gpt-4o-mini} is given a single one-shot
prompt and asked to write a unit test for one individual function of medium cyclomatic
complexity ($5 \le CC \le 10$) taken from a real open-source Java or Python project, does the
suite it produces achieve a median branch coverage of at least 80\,\% --- the level most
commonly quoted as an industry target --- measured over that function's own line range rather
than over the whole file?
\emph{H1: median branch coverage $\ge 80\,\%$.}

\textbf{RQ2-A --- Mutation threshold.} Does the same suite detect injected faults at a median
mutation score of at least 60\,\% --- a threshold set deliberately above the roughly 34\,\%
reported for LLM-generated suites in the studies reviewed in Section~\ref{sec:related} ---
measured with one operator set over the same line range?
\emph{H2-A: median mutation score $\ge 60\,\%$.}

\textbf{RQ2-B --- Comparison with baselines.} Matched function by function, does
\texttt{gpt-4o-mini} achieve a higher mutation score than the established automated generators
--- EvoSuite and Randoop for Java, Pynguin for Python --- when those tools are run on the same
targets under a fixed time budget?
\emph{H2-B: \texttt{gpt-4o-mini} scores above each baseline.}

\textbf{RQ3 --- Complexity and quality.} Within the narrow 5--10 band alone, does cyclomatic
complexity correlate negatively with the branch coverage the model achieves --- that is, does
test quality degrade measurably as complexity rises even inside a single band?
\emph{H3: Spearman $\rho < 0$.}

Table~\ref{tab:rqmap} maps each question to the test that answers it and to the data it is
answered on. Java and Python are reported separately throughout, because the toolchains and
the access rules of the two languages differ.

\begin{table}[t]
\centering
\caption{Each research question, its hypothesis, the statistical test used, and the sample it
is applied to. All tests are non-parametric at $\alpha = 0.05$.}
\label{tab:rqmap}
\begin{tabular}{llll}
\hline
RQ & Hypothesis & Test & Sample \\
\hline
RQ1   & median coverage $\ge 80\%$ & one-sample Wilcoxon vs.\ 80 & 60/lang.\ (all), effective \\
RQ2-A & median mutation $\ge 60\%$ & one-sample Wilcoxon vs.\ 60 & 60/lang.\ (all), measurable \\
RQ2-B & GPT $>$ each baseline      & paired Wilcoxon, by \texttt{function\_id} & matched pairs \\
RQ3   & Spearman $\rho < 0$        & Spearman correlation of CC vs.\ coverage & effective subset \\
\hline
\end{tabular}
\end{table}

\subsection{Dataset}
The proposal mentioned Defects4J~\cite{just2014defects4j} and CodeXGLUE as sources. We excluded the latter:
CodeXGLUE is a benchmark of code-intelligence tasks, including summarisation and
translation, not a set of functions amenable to coverage or mutation analyses; nor did
any of the papers surveyed for this proposal actually evaluate these tasks on CodeXGLUE
alongside any of these measurements. We deviated from the proposed protocol in that we mined
our own dataset instead of using Defects4J and CodeXGLUE directly; this change was documented
and approved by the instructor on 2026-07-02 (amendment~v1.2), before the experiment was run.
We therefore record functions from our own ten repositories at particular commits. Eight are Java projects overlapping
with Defects4J's subjects in their target programs: \texttt{commons-cli},
\texttt{commons-math}, \texttt{commons-csv}, \texttt{commons-collections}, \texttt{gson},
\texttt{jsoup}, \texttt{joda-time}, and \texttt{jfreechart}; two are Python projects,
\texttt{flask} and \texttt{requests}. URLs, licences, commit hashes, and dates of
download are recorded for every repository to facilitate reproduction; the final sample
contains 120 functions --- 60 in each language --- with cyclomatic complexity between 5
and 10 according to Lizard. Of these, 48 Python functions come from a single project,
\texttt{flask}, rather than being split evenly between two as the porting procedures
suggested; we retained this skew (resampling the dataset after observing the result
would itself have been a worse methodological choice) and discuss it at length in
Section~\ref{sec:threats}.

\subsection{Test generation}
\textbf{LLM.} All 120 test suites were generated by a single call to
\texttt{gpt-4o-mini-2024-07-18} (i.e., the dated snapshot; hence a rerun of this
experiment would use the same weights), \texttt{temperature=0}, \texttt{top\_p=1},
\texttt{max\_tokens=2048}, one-shot, with a worked example preceding the target function
in each prompt. We deviated from the proposed protocol in that we replaced \texttt{gpt-4o}
with \texttt{gpt-4o-mini-2024-07-18}: the team's API key had access to the smaller model only
at the moment the experiment was run. This change was documented on 2026-06-28 and approved by
the instructor on 2026-07-02 (amendment~v1.1), before the experiment was run. The replacement
model has been evaluated on the same task in this literature
before~\cite{konstantinou2026evaluating}, which keeps us within the evidential basis of this
report. None of the research questions, metrics, thresholds, or statistical tests were
affected by either amendment.

\textbf{Prompt.} Figure~\ref{fig:prompt} reproduces the Java prompt verbatim; the Python
prompt is identical except that it asks for \texttt{pytest} tests and names a module instead
of a package. The only per-function variation is the substitution of the target's qualified
name, invocation target, signature, callable kind, visibility, fixture requirement, package,
and source code.

\begin{figure}[t]
\begin{lstlisting}
You are an expert software tester.
Generate JUnit 5 tests to achieve maximum branch coverage and mutation score.
Do not output any explanations, only the executable Java code.

### Example Input ###
Class: org.example.MathUtils
Function: add
Code:
```java
package org.example;
public class MathUtils {
    public static int add(int a, int b) {
        if (a < 0 && b < 0) return 0;
        return a + b;
    }
}
```

### Example Output ###
```java
package org.example;
import org.junit.jupiter.api.Test;
import static org.junit.jupiter.api.Assertions.*;

public class MathUtilsTest {
    @Test
    void testAdd_BothNegative() { assertEquals(0, MathUtils.add(-1, -1)); }
    @Test
    void testAdd_Normal() { assertEquals(5, MathUtils.add(2, 3)); }
}
```

### Actual Task ###
Qualified target: `{qualified_name}`
Invocation target: `{invocation_target}`
Exact signature: `{signature}`
Callable kind: `{callable_kind}`; visibility: `{visibility}`.
Fixture requirement: `{fixture_required}`.
Generate JUnit 5 tests in package `{package_name}` that exercise exactly this target.
Include the correct `package {package_name};` at the top of the file.

```java
{source_code}
```
\end{lstlisting}
\caption{The one-shot prompt template, reproduced verbatim. Braces mark the fields
substituted per function.}
\label{fig:prompt}
\end{figure}

\textbf{Pipeline and edge cases.} Each function passes through five steps: (1) selection by
Lizard for $5 \le CC \le 10$; (2) prompt construction from the function's source and
metadata; (3) a single API call; (4) extraction of the fenced code block and writing of the
test file; (5) compilation, execution, and metric computation inside the original project.
Every call is a \emph{single attempt with no retry}: the call is wrapped in an exception
handler that, on any failure, records \texttt{gen\_status = api\_error} with an empty test
path, and the function is then scored \textsc{invalid} (0\%) exactly as a suite that failed
to compile. A response truncated at \texttt{max\_tokens} is not treated specially; it simply
fails to compile and is scored the same way. In the run reported here no LLM call failed:
all 120 functions returned \texttt{gen\_status = ok}. Timeouts did occur for the baselines
and are recorded as such --- three Randoop classes timed out, and the generation log
distinguishes \texttt{ok}, \texttt{failed}, and \texttt{timeout} for every tool invocation.

\textbf{Baselines.} EvoSuite~\cite{fraser2011evosuite} and
Randoop~\cite{pacheco2007randoop} generate the Java comparison suites; 60 seconds are
allocated to each class under test for test generation. That allocation was not fixed
initially: the pilot experiment spent 60s on EvoSuite and 10s on Randoop (a bias we only
became aware of when comparing the results), while the final run equalised the budgets
(see Section~\ref{sec:threats}). Pynguin~\cite{lukasczyk2022pynguin} is used for Python
with 90s per module under DYNAMOSA. Hypothesis was a viable baseline for the initial
comparison but was not actually used for the reported results.

\textbf{Replication package.} Our replication package, including the mined dataset with full
provenance, the exact prompts, the generation and measurement scripts, and the raw
per-function results, is publicly available at
\url{https://github.com/Nguyen-Tung-An/RT-SWT-nhom2}.

\subsection{Measurement}
We declined to measure functions in isolation; all test suites are executed in the target
project at the commit corresponding to the function being tested (recorded in the
dataset's provenance record; see the URLs, licences, hashes, and dates in the Dataset
subsection), with whatever environment and dependencies were present there at the time.
The reason for this design choice was that many of our sample functions only manifested
their documented behaviour while resolving their imported names inside the module or
package they belonged to (see the pilot result leading to this insight in
Section~\ref{sec:threats}).

\textbf{Java.} Tests are compiled and run with Maven~3.9.16; JaCoCo~0.8.12 measures branch
coverage while PIT~1.17.2~\cite{coles2016pit} performs mutation analysis. Both report
coverage and mutation scores at the class level, while our unit of analysis is the
function; hence the values reported for each function are the intersection of the
reported lines with the \texttt{[start\_line, end\_line]} of the target function across
all reports. The mutation score is the fraction of PIT mutants reported within that line
range to be killed by the test suite. There were two bugs in this pipeline affecting our
results: the PIT timeout was not handled by the measurement script, causing one entire
run to be dropped; and \texttt{jacoco.xml} and \texttt{mutations.xml} reports from the
previous function erroneously carried over into the next run in the same module. The
latter bug would have caused a function to erroneously report coverage of lines belonging
to another function (and hence artificially inflate its score), so the final script
clears the report directory between runs. Finally, GPT test suites had a much higher
fraction of compilation errors than the baseline tools'; hence we compiled and measured
each GPT test suite individually, per function, per run, rather than in bulk. This avoids
an entire test suite in a batch failing to compile because another test suite in the same
file has syntax errors.

\textbf{Python.} Functions are measured in an editable install of the pinned
\texttt{flask} (3.2.0.dev, pinned 2026-05-31) and \texttt{requests} (2.34.2, pinned
2026-06-15) checkouts; \texttt{coverage.py}~7.15.1 measures branch coverage, cut to the
function's line range as in Java. Mutation analysis is performed with two different
instruments, one for each test-suite flavour. The \texttt{gpt-4o-mini} suites are scored
with \texttt{mutmut}~2.4.4, with all lines outside the target function's range marked as
\texttt{\# pragma: no mutate}. Only test cases passing on the unmutated source are
considered; the score is the fraction of mutants with a non-surviving status (killed or
timed out) over all mutants generated in the line range. The Pynguin suites are scored
with a home-grown script using Python's standard-library \texttt{ast}: it replaces
\texttt{+} and \texttt{-}, \texttt{*} and \texttt{/}, \texttt{>} and \texttt{>=},
\texttt{<} and \texttt{<=}, \texttt{==}, \texttt{and}, and \texttt{or} with their opposites,
negates boolean literals,
increments numeric literals by 1, and applies any of these changes to exactly one site in
the target function's line range, up to 20 mutants per function. Every such mutant is
scored by rerunning the whole test suite with a 90s timeout: the suite killed the mutant
if it failed or timed out; the source code is reverted to its original state after each
run. Functions with no mutable site in range are reported as missing rather than scored 0,
since no mutants could be generated for them. Neither script attempts to detect equivalent
mutants; hence both report scores as a fraction of all mutants generated in range, for the
most conservative estimate. Both instruments use the same green-on-original gate and
line-range scoping; hence the only difference affecting the comparison is the set of
operators replaced. This makes the RQ2-B comparison between \texttt{gpt-4o-mini} and
Pynguin unsuitable for directly assessing the equivalence of their mutation scores; we
discuss this as a construct-validity threat in Section~\ref{sec:threats}. The Java RQ2-B
comparison is unaffected: PIT scores both test-suite flavours with the same instrument.
Limiting mutation analysis to the target-function line range is itself another
conservative choice: by avoiding helper functions and static variables, it prevents
incidental coverage and mutation of irrelevant code, maximising the score difference
between test-suite flavours for each function, at the expense of potentially missing
faults in wider contexts. Neither step involved any randomness, given the pinned checkout
and the model snapshot at \texttt{temperature=0}.

Any suite failing to compile, import, or run is marked \textsc{invalid}, with both metrics
scored at 0\% (see \S5.1 of the proposal). Before attempting to run the mutation analysis,
each test suite must pass on the original, unmutated source code. Mutation analysis run on
a failing test suite would artificially inflate the score: Section~\ref{sec:threats}
describes the pilot bug that led to this check. An \textsc{invalid} suite --- one that
failed to compile or run, or failed on the original source --- is not the same as a suite
that compiled and ran but failed to cover any of the target lines (i.e.,
\texttt{compiled=1}, coverage $=0$): both receive 0\% coverage, but for different reasons,
and both are tracked as separate outcome types in Section~\ref{sec:results} for easier
downstream analysis.

\subsection{Statistical analysis}
All tests are performed at $\alpha = 0.05$; all are non-parametric. We choose non-parametric
tests because coverage and mutation scores are percentages bounded between 0 and 100 and are
heavily zero-inflated in this sample, so normality was never assumed; this follows the
guidance of Arcuri and Briand for assessing randomized algorithms in software
engineering~\cite{arcuri2014hitchhiker}. RQ1 requires one test: a one-sample Wilcoxon
signed-rank test of branch coverage against the 80\% threshold. RQ2 requires two: a
one-sample Wilcoxon signed-rank test of mutation scores against the 60\% threshold, and a
paired Wilcoxon test (matched by \texttt{function\_id}) of the same against each baseline,
with matched-pairs rank-biserial $r$ as the effect size. Finally, RQ3 requires a Spearman
correlation of CC with branch coverage.

\textbf{Effect size and its interpretation.} Every $p$-value is reported with an effect size,
because significance alone says nothing about practical magnitude. For the Wilcoxon tests we
report the matched-pairs rank-biserial correlation $r$, and for RQ3 the Spearman $\rho$. We
interpret both on the same conventional bands: $|r| < 0.1$ negligible, $0.1 \le |r| < 0.3$
small, $0.3 \le |r| < 0.5$ medium, and $|r| \ge 0.5$ large. Because we run twelve tests in
total, we also report in Section~\ref{sec:threats} which conclusions survive a Bonferroni
correction to $\alpha' = 0.0042$. Every statistic is reported at two levels: \emph{all} ($N = 60$ per
language, with \textsc{invalid} or non-touching test suites scored at 0\% per the
pre-registered procedure) and \emph{effective} (test suites that successfully compiled and
touched the target lines). The two levels are never reported separately without the other,
since omitting the all-level would artificially inflate the statistics: an effective suite
has higher coverage than an invalid one (0\% versus some positive number), so the
effective figure is never lower than the all figure by construction (strictly higher for
coverage; for mutation the two can coincide at 0\%).
Section~\ref{sec:discussion} contains more context on this point.
